\documentclass[11pt,a4paper]{article}
\usepackage[utf8]{inputenc}
\usepackage[T1]{fontenc}
\usepackage{lmodern}
\usepackage[margin=2cm]{geometry}
\usepackage{listings}
\usepackage{xcolor}
\usepackage{titlesec}
\usepackage{booktabs}
\usepackage{array}
\usepackage{enumitem}
\usepackage{fancyhdr}
\usepackage{mdframed}
\usepackage{amsmath}
\usepackage{hyperref}

% ── Institutional colors ─────────────────────────────────────────────────────
\definecolor{salesiana}{RGB}{0,84,166}
\definecolor{cppbg}{RGB}{245,245,245}
\definecolor{kw}{RGB}{0,0,180}
\definecolor{cm}{RGB}{80,120,80}
\definecolor{str}{RGB}{160,0,0}
\definecolor{rubgreen}{RGB}{0,120,60}

\hypersetup{colorlinks=true, linkcolor=salesiana, urlcolor=salesiana}

% ── Section colors ────────────────────────────────────────────────────────────
\titleformat{\section}{\large\bfseries\color{salesiana}}{\thesection.}{0.5em}{}
\titleformat{\subsection}{\normalsize\bfseries\color{salesiana!80!black}}{\thesubsection.}{0.5em}{}

% ── Code style ───────────────────────────────────────────────────────────────
\lstset{
  backgroundcolor=\color{cppbg},
  basicstyle=\ttfamily\footnotesize,
  keywordstyle=\color{kw}\bfseries,
  commentstyle=\color{cm}\itshape,
  stringstyle=\color{str},
  breaklines=true,
  frame=single,
  rulecolor=\color{gray!40},
  numbers=left,
  numberstyle=\tiny\color{gray},
  numbersep=6pt,
  tabsize=4,
  showstringspaces=false,
  language=C++
}

% ── Header/footer ────────────────────────────────────────────────────────────
\pagestyle{fancy}
\fancyhf{}
\lhead{\small\textbf{Defense Guide --- Stereo Vision System}}
\rhead{\small Michael Lata $\cdot$ Jordy Romero --- G8}
\cfoot{\small\thepage}

% ── Criterion boxes ──────────────────────────────────────────────────────────
\newmdenv[linecolor=rubgreen,linewidth=1.5pt,
          frametitle={\color{rubgreen}\bfseries},
          frametitlebelowskip=4pt,
          innertopmargin=8pt]{critbox}

% =============================================================================
\begin{document}

% ── COVER PAGE (Spanish, institutional) ──────────────────────────────────────
\begin{titlepage}
\centering
\vspace*{0.8cm}
{\color{salesiana}\rule{\linewidth}{2pt}}\\[0.35cm]
{\Large\bfseries UNIVERSIDAD POLIT\'ECNICA SALESIANA}\\[0.15cm]
{\normalsize CARRERA DE INGENIER\'IA EN CIENCIAS DE LA COMPUTACI\'ON}\\[0.2cm]
{\color{salesiana}\rule{\linewidth}{1pt}}

\vspace{2.0cm}

\parbox{0.80\textwidth}{\centering\LARGE\bfseries\color{salesiana}
  Sistema de Estimaci\'on de Profundidad Estereoscp\'opica
  en Tiempo Real con Realidad Aumentada\\[0.3cm]
  \large mediante Microcontroladores ESP32-CAM
}

\vspace{1.4cm}

{\large\bfseries Gu\'ia de Defensa --- Visi\'on Artificial (G8)}\\[0.2cm]
{\normalsize Proyecto Integrador Final --- Parte I}

\vspace{1.6cm}

\begin{tabular}{@{}ll@{}}
\textbf{Autores:}    & Michael Lata $\cdot$ Jordy Romero \\[0.25cm]
\textbf{Docente:}    & Vladimir Robles Bykbaev \\[0.25cm]
\textbf{Asignatura:} & Visi\'on Artificial \\[0.25cm]
\textbf{Per\'iodo:}  & Abril -- Agosto 2026 \\[0.25cm]
\textbf{Fecha:}      & Junio 2026 \\
\end{tabular}

\vfill
{\color{salesiana}\rule{\linewidth}{2pt}}
\end{titlepage}

% ── Rubric overview ──────────────────────────────────────────────────────────
\begin{center}
{\large\bfseries\color{salesiana} DEFENSE GUIDE --- INTEGRATING PROJECT (Part I)}\\[6pt]
{\normalsize Real-Time Stereoscopic Depth Estimation System with Augmented Reality}
\end{center}

\vspace{4pt}
\begin{center}
\begin{tabular}{lcc}
\toprule
\textbf{Criterion} & \textbf{Weight} & \textbf{Main file(s)} \\
\midrule
1. Capture Architecture \& Hardware & 20\% & \texttt{ESP32\_CAM\_Stereo.ino, CameraSource.cpp} \\
2. Disparity Map Quality            & 25\% & \texttt{StereoProcessor.cpp} \\
3. Distance Precision \& Stability  & 20\% & \texttt{KalmanDepth.cpp, StereoProcessor.cpp} \\
4. Augmented Reality (AR) Effect    & 15\% & \texttt{ARRenderer.cpp} \\
5. Individual Oral Defense          & 20\% & (key questions) \\
\midrule
\textbf{Total} & \textbf{100\%} & \\
\bottomrule
\end{tabular}
\end{center}

\vspace{4pt}
\noindent\textbf{Note:} the 85\% weighting applies with ESP32-CAM-MB/XIAO-SP cameras;
this project uses ESP32-CAM AI-Thinker (\textbf{20\% weighting}).

% =============================================================================
\section{Criterion 1 --- Capture Architecture \& Hardware (20\%)}
% =============================================================================

\subsection*{What the rubric requires}
System built in OpenCV/C++ and communicating with ESP32-CAM microcontrollers whose
firmware parameters (AEC, AGC, AWB, XCLK) were deliberately tuned. C++ uses
\textbf{concurrent threads} for reading both cameras, reducing latency and network jitter.

\subsection*{ESP32-CAM Firmware --- \texttt{ESP32\_CAM\_Stereo.ino}}

\begin{critbox}[frametitle={Why AEC / AGC / AWB = OFF}]
SGBM computes disparity by comparing the intensity of corresponding pixels in the
left and right frames. If each camera adjusts its own exposure (AEC) independently,
the same physical point produces different pixel values in L and R, introducing
systematic errors in the disparity map (ghosting, black holes).
With the \textbf{same fixed exposure on both cameras}, photometric parity is guaranteed.
\end{critbox}

\begin{lstlisting}[caption={Fixed firmware parameters (lines 52--58)}]
constexpr int  CAMERA_XCLK_HZ  = 6000000; // 6 MHz - reduces WiFi jitter
constexpr bool CAMERA_AEC      = false;   // AUTO-EXPOSURE OFF -> manual
constexpr bool CAMERA_AGC      = false;   // AUTO-GAIN     OFF -> manual
constexpr bool CAMERA_AWB      = false;   // AUTO-BALANCE  OFF -> fixed color
constexpr int  CAMERA_EXPOSURE = 400;     // 0-1200; tuned per environment
constexpr int  CAMERA_GAIN     = 0;       // 0-30; minimum = less noise
\end{lstlisting}

\begin{critbox}[frametitle={Why XCLK = 6 MHz instead of the default 20 MHz}]
At higher clock frequencies, the sensor draws more current and generates more heat,
increasing jitter in the WiFi transmission. At 6 MHz the ESP32 prioritizes packet
integrity over speed, reducing the frame de-synchronization between L and R that
would otherwise degrade triangulation.
\end{critbox}

Parameters can be updated \textbf{at runtime without re-flashing} via HTTP:
\texttt{GET /setexposure?v=400} and \texttt{GET /setgain?v=0}.

\subsection*{Concurrent threads --- \texttt{CameraSource.cpp}}

\begin{lstlisting}[caption={Producer thread for HTTP camera frames}]
// Constructor launches the worker thread immediately
HttpCameraSource::HttpCameraSource(std::string url, int timeoutMs)
    : url_(url), timeoutMs_(timeoutMs) {
    worker_ = std::thread(&HttpCameraSource::workerLoop, this);
}
// Worker continuously downloads JPEG from /capture
void HttpCameraSource::workerLoop() {
    while (running_.load()) {
        // ... libcurl downloads JPEG ...
        {
            std::lock_guard<std::mutex> lk(mutex_);
            latestFrame_ = decoded;   // protected shared write
            hasFrame_    = true;
        }
        framesReceived_.fetch_add(1);
    }
}
// Main loop reads without blocking
bool HttpCameraSource::read(cv::Mat& out) {
    std::lock_guard<std::mutex> lk(mutex_);
    if (!hasFrame_) return false;
    out = latestFrame_.clone();
    return true;
}
\end{lstlisting}

\begin{itemize}[leftmargin=*]
  \item Left and right cameras each have their own dedicated producer thread.
  \item The main (controller) thread consumes the latest frame without waiting.
  \item \texttt{mutex\_} protects the shared frame between producer and consumer.
  \item \texttt{std::atomic<bool> running\_} stops the thread cleanly without a data race.
\end{itemize}

% =============================================================================
\section{Criterion 2 --- Disparity Map Quality (25\%)}
% =============================================================================

\subsection*{What the rubric requires}
Dense map, defined edges, low noise. Justified application of preprocessing (CLAHE)
and postprocessing (WLS) filters coupled to SGBM.

\subsection*{Full pipeline --- \texttt{StereoProcessor.cpp}}

\[
\text{Raw frame}
\xrightarrow{\text{remap}} \text{Rectified}
\xrightarrow{\text{cvtColor}} \text{Grayscale}
\xrightarrow{\text{CLAHE}} \text{Enhanced contrast}
\xrightarrow{\text{SGBM}} \text{Q16.4 disparity}
\xrightarrow{\text{WLS}} \text{Filtered disparity}
\xrightarrow{\text{reproject}} \text{3D point cloud}
\]

\begin{lstlisting}[caption={Step 1: Rectification (lens distortion correction)}]
// remap applies the pre-computed calibration maps to remove lens
// distortion and align epipolar lines horizontally
cv::remap(rawLeft,  result.leftRect,  calib_->map1x, calib_->map1y, INTER_LINEAR);
cv::remap(rawRight, result.rightRect, calib_->map2x, calib_->map2y, INTER_LINEAR);
\end{lstlisting}

\begin{lstlisting}[caption={Step 3: CLAHE preprocessing (lines 88--92)}]
// clipLimit=2.0: caps the maximum amplification of the local histogram
// tileSize 8x8: divides the image into tiles for local equalization
clahe_ = createCLAHE(cfg.claheClipLimit, Size(8, 8));
clahe_->apply(grayL, grayL);
clahe_->apply(grayR, grayR);
\end{lstlisting}

\begin{critbox}[frametitle={Why CLAHE before SGBM}]
SGBM finds correspondences by comparing block intensity. In areas with uneven
lighting or shadows, local contrast is low and the left block does not look like its
right counterpart. CLAHE equalizes the histogram \textbf{locally} without saturating
bright areas, improving discrimination at edges and fine textures.
\end{critbox}

\begin{lstlisting}[caption={Step 4: SGBM -- key parameters (lines 29--35)}]
sgbm_ = StereoSGBM::create(cfg.minDisp, numDisp, blockSz);
// P1 and P2 control the smoothness of the disparity map.
// The ratio P2 = 4*P1 follows Hirschmuller's recommendation.
sgbm_->setP1(8  * 3 * blockSz * blockSz);
sgbm_->setP2(32 * 3 * blockSz * blockSz);
sgbm_->setUniquenessRatio(cfg.uniquenessRatio);
sgbm_->setMode(StereoSGBM::MODE_SGBM_3WAY);
\end{lstlisting}

\begin{lstlisting}[caption={Step 5: WLS postprocessing (lines 101--108)}]
// Right matcher computes R->L disparity required by WLS
rightMatcher_->compute(grayR, grayL, disp16R);
Mat filtered;
// WLS fills holes and smooths edges (guided filter principle)
wlsFilter_->filter(disp16L, result.leftRect, filtered, disp16R);
filtered.convertTo(result.disparity, CV_32F, 1.0 / 16.0);
\end{lstlisting}

\begin{critbox}[frametitle={SGBM vs BM --- why we use SGBM}]
BM (Block Matching) searches only horizontally. SGBM (Semi-Global Block Matching)
adds smoothness constraints in 8 directions (semi-global), producing denser maps
with fewer holes, especially in low-texture regions.
\end{critbox}

\begin{lstlisting}[caption={Step 6: Invalid-zone fill (horizontal propagation)}]
// SGBM leaves ~numDisp black columns on the left with no valid data.
// The first valid value in each row is propagated leftward.
for (int y = 0; y < result.disparity.rows; ++y) {
    float* row = result.disparity.ptr<float>(y);
    float fill = 0.0f;
    for (int x = 0; x < result.disparity.cols; ++x) {
        if (row[x] > 0.5f) { fill = row[x]; break; }
    }
    if (fill <= 0.0f) continue;
    for (int x = 0; x < result.disparity.cols; ++x) {
        if (row[x] > 0.5f) break;
        row[x] = fill;
    }
}
\end{lstlisting}

\noindent\textbf{Modification note:} to fill \textbf{vertically} (column by column) instead,
make the outer loop iterate over \texttt{x} and the inner over \texttt{y}, switching from
\texttt{ptr<float>(y)[x]} (fast row-major access) to \texttt{at<float>(y,x)} (random access).

% =============================================================================
\section{Criterion 3 --- Distance Precision \& Stability (20\%)}
% =============================================================================

\subsection*{What the rubric requires}
Correct reprojection of the 3D matrix, Z coordinate in real units (cm/mm),
stable reading with a 1D Kalman Filter, fluctuation $\leq \pm 3\,\text{cm}$ in a static scene.

\subsection*{3D Reprojection}

\begin{lstlisting}[caption={Step 7: reprojectImageTo3D (line 131)}]
// Q is the reprojection matrix computed during stereo calibration.
// It converts the disparity map into a 3D point cloud (X, Y, Z in meters).
reprojectImageTo3D(result.disparity, result.points3D, calib_->Q, true);
\end{lstlisting}

\begin{lstlisting}[caption={Robust measurement: median over a 10x10 px region}]
float StereoProcessor::medianDepthAt(const Mat& points3D, int cx, int cy, int r) {
    // Collect all valid Z values in a square of radius r
    for (int y = y0; y <= y1; ++y)
        for (int x = x0; x <= x1; ++x) {
            const float z = points3D.at<Vec3f>(y, x)[2];
            // Keep only finite values in the plausible range (5 cm - 10 m)
            if (std::isfinite(z) && z > 0.05f && z < 10.0f)
                zVals.push_back(z);
        }
    // nth_element: O(n) - equivalent to median without sorting the entire vector
    std::nth_element(zVals.begin(), zVals.begin() + mid, zVals.end());
    return zVals[mid];
}
\end{lstlisting}

\subsection*{1D Kalman Filter --- \texttt{KalmanDepth.cpp}}

\begin{lstlisting}[caption={Scalar Kalman filter (the 3 key equations)}]
double KalmanDepth::update(double measurement) {
    // PREDICTION: uncertainty grows over time (process noise q)
    covariance_ += q_;

    // KALMAN GAIN: balances trust in prediction vs. new measurement
    const double gain = covariance_ / (covariance_ + r_);

    // UPDATE: corrects the estimate with the new measurement
    estimate_   += gain * (measurement - estimate_);
    covariance_  = (1.0 - gain) * covariance_;
    return estimate_;
}
\end{lstlisting}

\begin{center}
\begin{tabular}{lll}
\toprule
\textbf{Variable} & \textbf{Meaning} & \textbf{Effect if increased} \\
\midrule
\texttt{q\_}   & Process noise     & Kalman trusts measurement more (fast response) \\
\texttt{r\_}   & Measurement noise & Kalman trusts prediction more (smoother output) \\
\texttt{gain}  & Kalman gain       & 0 = ignore measurement; 1 = ignore prediction \\
\bottomrule
\end{tabular}
\end{center}

\noindent Key \textbf{C} resets the central Kalman filter (\texttt{kalmanCenter\_.reset()}).
Key \textbf{R} starts/stops CSV logging with columns: \texttt{frame, elapsed\_ms, z\_raw\_cm, z\_kalman\_cm, fps}.

\subsection*{Validation at 30 / 60 / 90 cm}
\texttt{python/collect\_measurements.py} automates data collection:
50 samples per distance, computes mean, standard deviation, and absolute error,
and verifies that \texttt{std} $\leq 3.0$ cm (rubric requirement).

% =============================================================================
\section{Criterion 4 --- Augmented Reality Effect (15\%)}
% =============================================================================

\subsection*{What the rubric requires}
An original computational visual effect overlaid on the rectified color transmission.
The effect must react \textbf{dynamically} to the measured distance (scale, position,
or behavior). Exclusive design unique to the group.

\subsection*{Effect description --- \texttt{ARRenderer.cpp}}

The \textbf{Energy Ball} effect has three stacked layers, all modulated by $Z$:

\subsubsection*{Layer 0: Depth-driven pixel distortion field (\texttt{applyDepthField})}

\begin{lstlisting}[caption={Depth-reactive wave distortion}]
// Strength inversely proportional to Z^3: drops off very fast with distance
const float z3       = zFocusMeters * zFocusMeters * zFocusMeters;
const float strength = std::clamp(0.6f / std::max(z3, 1e-4f), 0.0f, 70.0f);
const float speed    = 10.0f / std::max(0.15f, zFocusMeters);

// Sinusoidal pixel displacement for foreground pixels only
const float ox = strength * std::sin(y * 0.10f + t * speed);
const float oy = strength * std::cos(x * 0.10f + t * speed * 0.75f);
dst[x] = src.at<cv::Vec3b>(sy, sx);  // reads displaced pixel
\end{lstlisting}

\noindent\textbf{Dynamic behavior:} at 20 cm the wave displaces up to 70 px at high speed.
At 80 cm the perturbation drops to 2 px. Beyond 80 cm there is no effect (\texttt{strength < 2.0f}).

\subsubsection*{Layer 1: Radial color tint (\texttt{applyPixelVignette})}

\begin{lstlisting}[caption={Tint color reactive to distance}]
// Tint color changes with distance (in BGR):
if (depthCm < 40.0f)       { tB=20;  tG=60;  tR=255; } // RED  (danger: very close)
else if (depthCm < 80.0f)  { tB=20;  tG=220; tR=80;  } // GREEN (optimal zone)
else                        { tB=240; tG=160; tR=20;  } // CYAN  (far)
\end{lstlisting}

\subsubsection*{Layer 2: Pulsing aura with orbital sparks (\texttt{drawEnergyAura})}

\begin{lstlisting}[caption={Pulse and spark speed reactive to Z}]
// Soft 3 Hz pulse -- radius oscillates slightly
const float pulse  = 0.90f + 0.10f * std::sin(phase * 3.0);
// Closer -> more sparks and faster orbit
const float proximity  = std::clamp(20.0f / depthCm, 0.15f, 1.0f);
const int   sparkCount = 6 + int(8.0f * proximity);   // 6 to 14 sparks
const double spinSpeed = 1.0 + 2.0 * proximity;        // up to 3x faster
\end{lstlisting}

\begin{center}
\begin{tabular}{lll}
\toprule
\textbf{Distance} & \textbf{Color} & \textbf{Behavior} \\
\midrule
$<$ 40 cm (close)   & Red   & Strong wave, 14 fast sparks, intense tint \\
40--80 cm (optimal) & Green & Moderate wave, 10 sparks, soft tint \\
$>$ 80 cm (far)     & Cyan  & Minimal wave, 6 slow sparks \\
\bottomrule
\end{tabular}
\end{center}

% =============================================================================
\section{Criterion 5 --- Individual Oral Defense (20\%)}
% =============================================================================

\subsection*{Likely questions and key answers}

\textbf{1. What is matrix Q and what does \texttt{reprojectImageTo3D} do?}

$Q$ is the $4\times4$ reprojection matrix computed by \texttt{cv::stereoRectify}.
It transforms disparity $d$ into 3D coordinates:
\[
\begin{bmatrix} X \\ Y \\ Z \\ W \end{bmatrix} =
Q \cdot \begin{bmatrix} u \\ v \\ d \\ 1 \end{bmatrix}
\quad\Rightarrow\quad Z = \frac{f \cdot B}{d}
\]
where $f$ is the focal length and $B$ is the baseline (distance between cameras).
\texttt{reprojectImageTo3D} applies this operation to every pixel of the disparity map,
producing the 3D point cloud stored in \texttt{result.points3D}.

\vspace{6pt}
\textbf{2. How do threads work in the system?}

Each \texttt{HttpCameraSource} launches a \texttt{std::thread} (producer thread) that
continuously downloads JPEG frames via libcurl. The main thread (consumer) calls
\texttt{read()}, which grabs the latest frame under a \texttt{mutex}. This
producer-consumer architecture ensures that HTTP downloads do not block SGBM processing.

\vspace{6pt}
\textbf{3. How are SGBM parameters tuned?}

\begin{itemize}[leftmargin=*,noitemsep]
  \item \texttt{numDisparities}: search range in pixels (must be a multiple of 16). Larger = sees farther but slower.
  \item \texttt{blockSize}: comparison block size (odd, $\geq 3$). Larger = smoother but loses detail.
  \item \texttt{P1, P2}: smoothness penalties. $P2 = 4 \cdot P1$ follows Hirschmuller's guideline. High $P2$ = smooth plane transitions.
  \item \texttt{uniquenessRatio}: rejects a match if the second-best candidate is within X\% of the best.
  \item \texttt{speckleWindowSize / Range}: removes isolated blobs (salt-and-pepper noise).
\end{itemize}

All parameters are adjustable \textbf{in real time} via the trackbar control window --- no restart needed.

\vspace{6pt}
\textbf{4. What does key C do?}

It resets the central Kalman filter (\texttt{kalmanCenter\_.reset()}), setting
\texttt{initialized\_ = false}. The next measurement is accepted directly as the new
initial estimate, discarding all previous history. Useful when the target object
changes position abruptly.

\vspace{6pt}
\textbf{5. Why use the median instead of the mean to measure Z?}

The median is robust to outliers. Pixels at object edges have erroneous disparities
due to occlusion. The mean would be contaminated by those values; the median ignores
them as long as more than 50\% of pixels in the region are valid.

\vspace{6pt}
\textbf{6. What is the difference between Z raw and Z Kalman in the HUD?}

\texttt{Z raw} is the direct output of \texttt{medianDepthAt} (may vary frame to frame).
\texttt{Z Kalman} is the filtered estimate: more stable, with a slight lag, but meets
the $\leq \pm 3\,\text{cm}$ fluctuation requirement of the rubric.

\vspace{6pt}
\textbf{7. Why does the AR effect apply only to the foreground?}

In \texttt{applyDepthField} there is a threshold: \texttt{const float fgLimit = zFocusMeters + 0.50f}.
Only pixels with $Z < \text{focus} + 0.5\,\text{m}$ receive the distortion.
Background pixels (wall, floor) are untouched, giving the visual impression that
the effect wraps only around the measured object.

\vspace{6pt}
\textbf{8. If asked to change Step 6 fill direction from horizontal to vertical, what changes?}

The outer loop changes from \texttt{for y} to \texttt{for x}, and the inner from
\texttt{for x} to \texttt{for y}. Access changes from \texttt{ptr<float>(y)[x]}
(efficient row-major) to \texttt{at<float>(y,x)} (random access, required when
\texttt{x} is the outer axis).

% =============================================================================
\section*{File-to-criterion reference}
% =============================================================================

\begin{center}
\begin{tabular}{p{4.2cm}p{5.5cm}p{4.5cm}}
\toprule
\textbf{File} & \textbf{Responsibility} & \textbf{Criterion} \\
\midrule
\texttt{ESP32\_CAM\_Stereo.ino}  & Firmware: AEC/AGC/AWB, XCLK, HTTP endpoints & 1 \\
\texttt{CameraSource.cpp/.hpp}   & Producer thread, mutex, non-blocking read     & 1 \\
\texttt{StereoCalibration.cpp}   & Loads \texttt{stereo\_calib.yml}, rectification maps, Q & 2, 3 \\
\texttt{StereoProcessor.cpp}     & remap $\to$ CLAHE $\to$ SGBM $\to$ WLS $\to$ reproject & 2, 3 \\
\texttt{KalmanDepth.cpp}         & Scalar 1D Kalman filter                       & 3 \\
\texttt{ARRenderer.cpp}          & Distortion field, radial tint, pulsing aura   & 4 \\
\texttt{DisplayManager.cpp}      & HUD, trackbars, windows                       & 1, 2, 3 \\
\texttt{StereoController.cpp}    & Main MVC loop, CSV logger, key handling       & all \\
\texttt{collect\_measurements.py}& 30/60/90 cm validation, CSV results           & 3 \\
\texttt{find\_cameras.py}        & Network scan to locate ESP32-CAMs             & 1 \\
\bottomrule
\end{tabular}
\end{center}

\end{document}
